\documentclass{resume} % Use the custom resume.cls style

\usepackage[left=0.4 in,top=0.4in,right=0.4 in,bottom=0.4in]{geometry} % Document margins
\usepackage{xcolor}
\definecolor{softorange}{RGB}{255,127,0}
\newcommand{\tab}[1]{\hspace{.2667\textwidth}\rlap{#1}} 
\newcommand{\itab}[1]{\hspace{0em}\rlap{#1}}
\name{Wuxinlin Cheng} % Your name
% You can merge both of these into a single line, if you do not have a website.
%\address{+1(123) 456-7890 \\ San Francisco, CA} 
%\address{\href{mailto:wcheng7@stevens.edu}{wcheng7@stevens.edu} \\ C: 551-689-1091
\address{Email: \href{mailto:chengwu1230@gmail.com}{chengwu1230@gmail.com} 
\\ Web: \href{https://chengwuxinlin.github.io/}{chengwuxinlin.github.io} }  %

\begin{document}
%----------------------------------------------------------------------------------------
%----------------------------------------------------------------------------------------
%	EDUCATION SECTION
%----------------------------------------------------------------------------------------

\begin{rSection}{Education}
\item[]
Stevens Institute of Technology\textnormal{ Hoboken, NJ - Doctor of Philosophy in Computer Engineering \hfill {Jan 2021 – May 2025}}\\
Stevens Institute of Technology\textnormal{ Hoboken, NJ - Master of Engineering in Computer Engineering \hfill {Jan 2019 – Dec 2020}}\\
Sichuan University\textnormal{ Chengdu, China - Bachelor of Engineering in Electrical Engineering \hfill {Sep 2014 – Jun 2018}}
\end{rSection}
%----------------------------------------------------------------------------------------
% Projects
%----------------------------------------------------------------------------------------
\begin{rSection}{Projects}
\item \textbf{SPADE: A Metric for Black-Box Adversarial Robustness Evaluation (ICML 2021)}
\begin{itemize}
    \item Developed SPADE, a novel metric for evaluating adversarial robustness in ML models, employing bijective distance mappings of input/output graph-based manifolds. Enabled downstream applications such as adversarial training by revealing data robustness, resulting in up to $18\%$ enhanced accuracy compared to standard PGD training.
\end{itemize}

\item \textbf{Improving Out-of-Distribution (OOD) Robustness with Diffusion Model (KLA Corporation)}
\begin{itemize}
    \item Enhanced OOD robustness in a multi-label regression task using diffusion models. Developed a novel, brand-new diffusion model capable of conditionally generating extremely high-dimensional spectra OOD data. Achieved a 40\% improvement on in-distribution data and a 25\% improvement on OOD data.
\end{itemize}

\item \textbf{CirSTAG: Graph Neural Networks Stability Analysis \textcolor{softorange}{(Best Paper Nominee, top$\approx$0.5\%, DAC 2025)}}
\begin{itemize}
    \item Developed CirSTAG, a spectral framework for analyzing the stability of Graph Neural Networks (GNNs) using probabilistic graphical models. CirSTAG outperforms traditional Netattack in diverse GNNs, achieving 20\% higher error rates in adversarial attacks and reducing defense error rates by 50\%. Additionally, CirSTAG offers a significant improvement in evaluating and enhancing large-scale integrated circuit design. Moreover, \textbf{CirSTAG can identify user stability in recommendation systems.}
\end{itemize}

\item \textbf{AI Agent for Education, Entertainment, and Company in Virtual Reality }(Stealth Startup)
\begin{itemize}
    \item Led a stealth startup lab in developing 3D avatar assistants, establishing an innovative framework integrating Natural Language Processing and 3D model reconstruction. Developed a language emotion recognition system using Bert with 82\% accuracy across 27 emotions and pioneered a text-to-speech model based on VIST2, achieving a 4.4 Mean Opinion Score and a 0.014 Real-Time Factor. Collaborated on facial reconstruction, LLM dialogue, and real-time lip-sync technologies.
\end{itemize}


\item \textbf{Real-Time Talking Heads from User-Uploaded Avatars (CVPR 2024: \href{https://drive.google.com/file/d/1HsqRwAYIS7WWpbNsqjjuBPfROzYh0opo/view?usp=sharing}{Video Example})}
\begin{itemize}
    \item Developed a state-of-the-art system capable of generating real-time talking heads from user-uploaded avatar images, synchronized with any audio input. Pioneered the integration of this technology with LLM for live, avatar-based conversations, enhancing user interaction and engagement across various digital platforms. This innovation opens new interactive possibilities in social media, gaming, e-learning, and virtual customer support by facilitating more natural and engaging digital communication experiences.
\end{itemize}



\item \textbf{Model Accuracy \& Runtime Improvement for Vial Classification (LEMA (Beijing) Technology Co., Ltd.)}
\begin{itemize}
    \item Improved accuracy and reduced computational load for vial classification models by leveraging data selection, neural network pruning, and robust training strategies. Achievements include a significant $42\%$ improvement in model accuracy and a $5.8\times$ speedup in runtime over previous industrial models.
\end{itemize}

\item \textbf{Hair Detection Stability Improvement (LEMA (Beijing) Technology Co., Ltd.)}
\begin{itemize}
    \item Enhanced stability and performance of hair detection models through robust training and data augmentations applied to a small dataset (200 pictures). Achieved exceptional test accuracy of 99.97\% on the LEMA (Beijing) Technology Co., Ltd model. Notably, the model's accuracy remained consistent even under severity 5 corruption, demonstrating robustness to reasonable perturbations such as brightness differences or blur.
\end{itemize}


\end{rSection} 

\begin{rSection}{PERSONAL EXPERIENCE}
\item[]
AI Algorithm Engineer (Deep Learning)\textnormal{- KLA Corporation}\textnormal{\hfill {August 2024 – Present}} \\
Conference Reviewer \textnormal{- ICLR \hfill {2026}} \\
Research Assistant \textnormal{- Stevens Institute of Technology}\textnormal{\hfill {June 2020 – May 2025}} \\
AI Algorithm Engineer Intern (Deep Learning)\textnormal{- KLA Corporation}\textnormal{\hfill {May 2024 – August 2024}} \\
\textbf{Scientific Advisor/Lead of AI Lab} \textnormal{- Stealth Startup       \hfill {Sep 2023 – April 2024}} \\%AskAlora, Interest Group AI 
\textbf{Scientific Advisor} \textnormal{- LEMA (Beijing) Technology Co., Ltd.\hfill {Dec 2021 – Sep 2023}} \\
\textbf{Research Intern} \textnormal{- Shanghai ASES Spaceflight Technology Co. Ltd.\hfill {Sep 2018 – Jan 2019}} 

\end{rSection}

%----------------------------------------------------------------------------------------
% TECHINICAL STRENGTHS	
%----------------------------------------------------------------------------------------
\begin{rSection}{SKILLS}
\item[]
\begin{tabular}{ @{} >{\bfseries}l @{\hspace{6ex}} l }
Research Areas & \textnormal{Graph Learning, ML, CV, AI Stability, NLP, AIGC, Diffusion Model, LLM, GNN, Model Compression} \\
Software \& Tools & \textnormal{Python, C++, Pytorch, Tensorflow, Scikit-learn} \\
\end{tabular}\\
\end{rSection}


%----------------------------------------------------------------------------------------
\begin{rSection}{AWARDS AND HONORS} 
\item[]
\begin{itemize}
    \item 	\textnormal{Outstanding Ph.D. Dissertation Award in Computer Engineering, USA \hfill{May 2025}}
    \item 	\textnormal{Stevens Institute of Technology scholarship, USA \hfill{Jan 2020}}
    \item 	\textnormal{National Entrepreneurship (top $1\%$ nationally) at 3rd China International College “Internet Plus” Competition \hfill{Jan 2018}}
    \item 	\textnormal{Sichuan University scholarship, China \hfill{Dec 2017}}
    \item 	\textnormal{Honorable Award (top $21\%$ globally) at the Mathematical Contest in Modeling, USA \hfill{Apr 2017}}
\end{itemize}


\end{rSection}


\begin{rSection}{Publication}
\item \textbf{Wuxinlin Cheng}. \textnormal{``Machine-learning solution for Cu dishing metrology in advanced packaging.'' \textit{SPIE Advanced Lithography + Patterning}, \textcolor{softorange}{[Conference Presentation]}, Paper 13981--117, 2026.}


\item \textbf{Wuxinlin Cheng, \textnormal{Yupeng Cao, Jinwen Wu, Koduvayur Subbalakshmi, Tian Han, and Zhuo Feng. “SALMAN: Stability Analysis of  Language Models Through the Maps Between Graph-based Manifolds.” ICLR 2026 Under Review}}
\item \textbf{Wuxinlin Cheng, \textnormal{Chenhui Deng, Ali Aghdaei, Zhiru Zhang, and Zhuo Feng. “SAGMAN: Stability Analysis of Graph Neural Networks on the Manifolds.” ICLR 2026 Under Review}}
\item \textbf{Wuxinlin Cheng, \textnormal{Yihang Yuan, Chenhui Deng, Ali Aghdaei, Zhiru Zhang, and Zhuo Feng. “CirSTAG: Circuit Stability Analysis via Graph
Neural Networks.” Design Automation Conference (DAC), 2025}} 
\item \textbf{Wuxinlin Cheng, \textnormal{Yupeng Cao, Cheng Wan, Sihan Chen. “Towards a Real-Time Interactive Framework for Talking Avatars.” The IEEE/CVF Conference on Computer Vision and Pattern Recognition (CVPR) demo 2024}}
\item \textnormal{John Anticev, Ali Aghdaei, }\textbf{Wuxinlin Cheng, }\textnormal{and Zhuo Feng. “SGM-PINN: Sampling Graphical Models for Faster Training of Physics-Informed Neural Networks.” Design Automation Conference (DAC), 2024}
\item \textbf{Wuxinlin Cheng, \textnormal{Chenhui Deng, Zhiqiang Zhao, Yaohui Cai, Zhiru Zhang, and Zhuo Feng. “SPADE: A Spectral Method for Black-Box Adversarial Robustness Evaluation.” International Conference on Machine Learning (ICML), 2021.}} 
\item \textbf{Wuxinlin Cheng, \textnormal{Xu Zhou, Xin He. “Quick Pass Optimization in Airport Security Check,” TianFuShuXue, ISSN: 1006-0324, Vol, 21, 2018}}


\end{rSection}


\begin{rSection}{Intellectual Property}
\item \textbf{KLA Corporation}. \textnormal{\textcolor{softorange}{[Internal Invention Disclosure]} Developed a proprietary feature-encoding method in metrology modeling, improving both in-distribution accuracy and out-of-distribution robustness.}



\item \textbf{KLA Corporation}. \textnormal{\textcolor{softorange}{[Internal Invention Disclosure]} Proprietary generative modeling approach for synthesizing out-of-distribution metrology data to improve model robustness.}


\end{rSection} 

%----------------------------------------------------------------------------------------



\end{document}
